\documentclass[aspectratio=169,13pt,fleqn] {beamer}
%
% Choose how your presentation looks.
%
% For more themes, color themes and font themes, see:
% http://deic.uab.es/~iblanes/beamer_gallery/index_by_theme.html
%
\mode<presentation>
{
	\usetheme{Madrid}      % 整体主题
	\usecolortheme{beaver} % 配色主题（棕红系）
	\usefonttheme{serif}   % 全局衬线字体（正文宋体类）
	\setbeamertemplate{items}[default] % 【items父模板】控制 itemize+enumerate 符号
%	\setbeamertemplate{items}{\Large\beamertemplateballitem}
	\setbeamertemplate{sections/subsections in toc}[ball] % 目录里章节用圆圈标记
	
	%  \setbeamertemplate{navigation symbols}{} % 取消注释可隐藏底部导航小图标
	\setbeamertemplate{caption}[numbered] % 图/表标题带编号：图1、表1
}

\usepackage[,UTF8,linespread=1.2]{ctex}
\usepackage{wrapfig}
\usepackage{caption}
\captionsetup[figure]{name=图,labelsep=period}



\usepackage{xcolor}
\definecolor{xy_blue}{cmyk}{0.6,0.35,0.0,0.4}
\definecolor{xy_orange}{cmyk}{0.0,0.7,1.0,0.0}
\definecolor{xy_green}{cmyk}{0.2,0.0,1.0,0.15}
\xdefinecolor{xy_gold}{RGB}{135,103,79}
\xdefinecolor{xy2_gold}{cmyk}{0,0.35,0.75,0.05}

\setbeamercolor{section number projected}{fg=red,bg=red!20}%控制球的颜色
\setbeamercolor{subsection number projected}{fg=blue,bg=blue!40}
\setbeamercolor{item}{fg=xy_gold} %设置列表项颜色

\usepackage{amsthm}
\newtheorem{peroration}{结论：}
\newtheorem*{exploration}{探究：}
\newtheorem*{prooft}{证明：}
\newtheorem*{solve}{解:} 
\newtheorem{instance}{例}   % 解：独立计数器
     

% 配色（你前面需要的，每个环境单独颜色）
\setbeamercolor{theorem}{fg=black,bg=blue!10}
\setbeamercolor{solution}{fg=black,bg=green!10}

\BeforeBeginEnvironment{peroration}{%结论： 
	\setbeamercolor{block title}{bg=xy_green,fg=white}
	\setbeamercolor{block body}{bg=xy_green!10}
}
\AfterEndEnvironment{peroration}{
	\setbeamercolor{block title}{bg=xy_blue,fg=white}
	\setbeamercolor{block body}{bg=xy_blue!10}
}
\BeforeBeginEnvironment{exploration}{%探究
	\setbeamercolor{block title}{bg=xy_orange,fg=white}
	\setbeamercolor{block body}{bg=xy_orange!10}
}
\AfterEndEnvironment{peroration}{
	\setbeamercolor{block title}{bg=xy_blue,fg=white}
	\setbeamercolor{block body}{bg=xy_blue!10}
}
\BeforeBeginEnvironment{solution}{%解
	\begingroup
	\setbeamercolor{block title}{bg=xy_green,fg=white}
	\setbeamercolor{block body}{bg=xy_green!10}
}
\AfterEndEnvironment{solution}{
	\endgroup
}
\BeforeBeginEnvironment{prooft}{%证明
	\setbeamercolor{block title}{bg=xy_green,fg=white}
	\setbeamercolor{block body}{bg=xy_green!10}
}
\AfterEndEnvironment{prooft}{
	\setbeamercolor{block title}{bg=xy_green,fg=white}
	\setbeamercolor{block body}{bg=xy_green!10}
}
\BeforeBeginEnvironment{solve}{%解：
	\setbeamercolor{block title}{bg=xy_orange!60,fg=white}
	\setbeamercolor{block body}{bg=xy_orange!10}
}
\AfterEndEnvironment{solve}{%解：
	\setbeamercolor{block title}{bg=xy_orange!60,fg=white}
	\setbeamercolor{block body}{bg=xy_orange!10}
}
   
\BeforeBeginEnvironment{instance}{% 例子
	\begingroup
   	\setbeamercolor{block title}{bg=xy_gold!40,fg=black}
   	\setbeamercolor{block body}{bg=xy_gold!10}
   }
\AfterEndEnvironment{instance}{
	\endgroup
   }     
   
\renewcommand{\thefigure}{\arabic{figure}}

\title[2.2基本不等式]{2.2~~基本不等式(第2课时)}
\author[好好学习，努力向上]{}
\institute{}
\date{2026.9.23}
\setbeamertemplate{theorems}[numbered]

\AtBeginSection[]
{
  \begin{frame}<beamer>
    \frametitle{本节主要内容}
    \tableofcontents[currentsection,currentsubsection]
  \end{frame}
}

\begin{document}
\begin{frame}{}
	\titlepage
\end{frame}

\begin{frame}{内容概览}
	\tableofcontents
\end{frame}

\section{前情回顾}
\begin{frame}{前情回顾}
	
	\begin{enumerate}[(1)]
		\item 重要不等式：已知$a$、$b\in R$，则
		\[
		\alert{a^2 + b^2 \geqslant 2ab}
		\]
		当且仅当$a=b$时等号成立。
		\item 基本不等式：已知$a$，$b\in R_{\;+}$，则
		\[
		\alert{\sqrt{ab} \leqslant \frac{a+b}{2}}
		\]
		当且仅当$a=b$时等号成立。
	\end{enumerate}
	
\end{frame}

\section{基本不等式的应用}
\begin{frame}[t]{基本不等式的应用1}
	\begin{instance}
		已知 $x>1$，求 $x+\frac{1}{x-1}$的最小值。
	\end{instance}
	
	
\end{frame}

\begin{frame}[t]{基本不等式的应用2}
	\begin{instance}
		求$\sqrt{x\;(10-x)}$的最大值。
	\end{instance}
	
	
\end{frame}

\begin{frame}[t]{基本不等式的应用3}
	\setbeamertemplate{enumerate items}[default]
	\begin{instance}
		

	\begin{enumerate}[(1)]
		\item 已知 $x>0$，$y>0$，$x+y=1$，求$\dfrac{1}{x}+\dfrac{2}{y}$的最小值；
		\item 已知 $a>0$，$b>0$，且$\dfrac{1}{a}+\dfrac{9}{b} = 1$，求$a+b$的最小值。
	\end{enumerate}	
	\end{instance}
\end{frame}


\begin{frame}{答案与解析---}
	
	\begin{solve}[1]
	因为$x>0$，$y>0$，且$x+y=1$，所以
	\begin{align*}
		&\frac{1}{x}+\frac{2}{y} = \left(\frac{1}{x} + \frac{2}{y} \right)(x+y)\\
		&=1+\frac{2x}{y} + \frac{y}{x} +2\geqslant 3+2\sqrt{\frac{2x}{y}\cdot\frac{y}{x}}\\
		&=3+2\sqrt{2}
	\end{align*}
	当且仅当$\frac{2x}{y} = \frac{x}{y}$，且$x+y=1$时取等号，解得$x=\sqrt{2}-1,y=2-\sqrt{2}$。
	
	所以当$x=\sqrt{2}-1,y=2-\sqrt{2}$时，原式有最小值$3+2\sqrt{2}$。
	\end{solve}
\end{frame}

\begin{frame}{答案与解析---}
	
	\begin{solve}[2]
		因为$a>0$，$b>0$，且$\frac{1}{a} + \frac{9}{b} =1$，所以
		\begin{align*}
			& a+b= \left(\frac{1}{a} + \frac{9}{b} \right)(a+b)\\
			&=1+\frac{b}{a} + \frac{9a}{b} +9\geqslant 10+2\sqrt{\frac{b}{a}\cdot\frac{9a}{b}}\\
			&=10+2\times 3 =16
		\end{align*}
		当且仅当$\frac{b}{a} = \frac{9a}{b}$，且$\frac{1}{a}+\frac{9}{b}=1$时取等号，解得$a=4,b=12$。
		
		所以当$a=4,b=12$时，原式有最小值$16$。
	\end{solve}
\end{frame}



\begin{frame}{布置作业}

\end{frame}

\end{document}
